\begin{table*}[t]
\centering
\small
\caption{Counter-speech generated by Gemma-4-E4B for EDOS messages where the label predicted by the EDOS classifier differs from the gold annotation. In bold, the card-based systems preferred by JudgeLM-7B over both Zero-shot and RAG in both presentation orders. EDOS label: the label predicted by the EDOS classifier (used by the + pred. label systems) and the gold annotation (used by the + gold label systems).}
\label{tab:examples-edos-label-mismatch}
\begin{tabularx}{\textwidth}{@{}p{0.2\textwidth}X@{}}
\toprule
\textbf{Hate speech} & \textbackslash{}>woman here Opinion discarded. \\
\textbf{EDOS label (pred.)} & 3.2 immutable gender differences and gender stereotypes \\
\textbf{EDOS label (gold)} & 2.1 descriptive attacks \\
\midrule
\textbf{System} & \textbf{Generated counter-speech} \\
Zero-shot & Everyone deserves to have their opinions heard, regardless of gender. Let's commit to respectful dialogue where all voices are valued. \\
RAG & My opinions are valid, and as women, we are experts on our own experiences and what we need, so please respect my voice. Silencing women through attacks undermines our right to free speech and limits our participation in public discourse. \\
RAG + pred. label & Dismissing a woman's opinion based on gender reinforces harmful stereotypes that limit personal abilities and choices, which is a core aspect of sexism. All individuals, regardless of gender, deserve to have their voices heard and their perspectives valued. \\
RAG + gold label & Dismissing a woman's opinion is a form of sexist hate speech that aims to undervalue her skills and opinions, which is a severe form of gender-based violence. All individuals, regardless of gender, deserve to have their voices heard and their perspectives respected. \\
\textbf{Card-RAG} & \textbf{Dismissing a woman's opinion is a form of sexist hate speech that aims to devalue her skills and silence her voice. Every woman deserves to have her opinions respected without fear of being dismissed or pressured to change her behavior.} \\
\textbf{Card-RAG + pred. label} & \textbf{Dismissing a woman's opinion based on gender reinforces harmful stereotypes that limit everyone's ability to develop and contribute. Valuing diverse perspectives is crucial because sexist attitudes and biases contribute to exclusion and limit opportunities for all.} \\
\textbf{Card-RAG + gold label} & \textbf{Dismissing a woman's opinion is a form of sexist hate speech that undervalues her skills and voice, which undermines her freedom of speech. All individuals deserve to have their perspectives heard without being subjected to derogatory generalizations.} \\
Hybrid RAG & My opinions are valid, and dismissing them is a form of sexist hate speech that aims to undervalue women's skills and silence their voices. I deserve to participate in discussions without being told I shouldn't speak on certain topics. \\
Hybrid + pred. label & Dismissing someone's opinion based on gender reinforces harmful gender stereotypes, which are linked to sexism that limits opportunities and perpetuates discrimination for everyone. True equality requires valuing all voices, regardless of gender, because sexist attitudes contribute to systemic inequalities. \\
Hybrid + gold label & Every individual's opinion holds value, and dismissing a woman's voice contributes to the silencing and objectification that sexist hate speech aims to enforce. We must challenge these harmful generalizations because they undermine freedom and perpetuate severe psychological and emotional harm. \\
\midrule[\heavyrulewidth]
\textbf{Hate speech} & Just take it easy on the weights so you don't end up with the big biceps and linebacker shoulders too many women have these days. \\
\textbf{EDOS label (pred.)} & 3.2 immutable gender differences and gender stereotypes \\
\textbf{EDOS label (gold)} & 3.3 backhanded gendered compliments \\
\midrule
\textbf{System} & \textbf{Generated counter-speech} \\
Zero-shot & Fitness goals are personal, and strength is a positive attribute for everyone. Focusing on health and personal goals is what truly matters. \\
RAG & Sport should celebrate athletic achievement regardless of gender, as cultural life, including sport, is a key shaper of attitudes and gender roles (Evidence 4). Normalised notions about gender can prevent us from accepting that women can be champions just as much as men, and we must challenge these stereotypes (Evidence 1). \\
RAG + pred. label & Physical strength and body shape are not determined by gender stereotypes, which arbitrarily assign characteristics to men and women. Valuing certain physical traits over others reinforces harmful biases that limit personal development and choices for everyone. \\
RAG + gold label & Physical strength and body shape are not defined by gender stereotypes; limiting men's physical expression based on outdated norms reinforces harmful biases. Everyone has the right to pursue their physical goals without being judged by sexist expectations. \\
\textbf{Card-RAG} & \textbf{Body size and strength are not defined by gender stereotypes, which often police women's bodies and demand specific beauty standards. Everyone deserves to pursue their physical goals without facing judgment based on outdated gender norms.} \\
\textbf{Card-RAG + pred. label} & \textbf{Physical strength and body shape are not determined by gender, as there is no evidence that men's and women's brains are wired differently. Stereotypes about what men should look like limit individual potential and reinforce harmful expectations.} \\
\textbf{Card-RAG + gold label} & \textbf{Physical strength is not defined by gender, and equating women's strength with 'big biceps' reinforces harmful stereotypes that limit women's capacity to pursue physical abilities. We should focus on describing actions and qualities without gendered labels to avoid belittling or implying inferiority.} \\
Hybrid RAG & Body diversity and strength are not dictated by gender; framing women's bodies as objects or policing their appearance ignores the double standards that demand women meet specific beauty standards while labeling natural states as lewd. We should challenge these restrictive gender stereotypes to accept all bodies and abilities equally. \\
Hybrid + pred. label & Physical strength and body shape are not defined by gender stereotypes, which often limit men by associating them with specific traits like aggressiveness or leadership. Everyone has the right to pursue their personal abilities and make choices about their bodies without judgment based on outdated gender norms. \\
Hybrid + gold label & Physical strength and body shape are not defined by gender; equating certain physical traits with femininity or weakness is a sexist stereotype. Everyone has the right to pursue physical activities and develop their body in a way that feels right for them. \\
\midrule[\heavyrulewidth]
\textbf{Hate speech} & The View. Empty head emotional woman's daily bitchcraft lesson. \\
\textbf{EDOS label (pred.)} & 2.2 aggressive and emotive attacks \\
\textbf{EDOS label (gold)} & 2.1 descriptive attacks \\
\midrule
\textbf{System} & \textbf{Generated counter-speech} \\
Zero-shot & Everyone deserves to be treated with respect, regardless of their opinions or appearance. Engaging in personal attacks only detracts from meaningful conversation. \\
RAG & Instead of perpetuating stereotypes, media can powerfully showcase strong, intelligent, and diverse women as resilient leaders and activists, highlighting their strength rather than just victimisation. We must actively promote these empowering narratives to challenge harmful norms and foster a more equitable conversation. \\
RAG + pred. label & Reducing women to insults ignores the severe, real impacts of sexist hate speech, which aims to humiliate and control women's voices and participation. Every woman deserves to be treated with respect, not subjected to attacks based on gender. \\
RAG + gold label & Labeling someone as 'emotional' or 'empty-headed' is a derogatory attack that dismisses their views, which is a form of sexist hate speech that aims to undervalue women's opinions. Instead of resorting to insults, we should engage with ideas respectfully, recognizing that passionate expression is a valid human trait. \\
Card-RAG & Dismissing critical voices as 'emotional' ignores the real and severe psychological impacts of sexist hate speech, which aims to devalue women's opinions and silence them. Every woman and girl is a potential target, and challenging norms is a vital part of freedom of speech. \\
\textbf{Card-RAG + pred. label} & \textbf{Dismissing someone's views as mere 'bitchcraft' is a form of sexist hate speech that aims to devalue women's opinions and silence them, which undermines freedom of speech. Every woman has the right to express her views without facing insults or being objectified based on her gender.} \\
\textbf{Card-RAG + gold label} & \textbf{Labeling women with derogatory terms like 'emotional' or 'empty head' is a form of descriptive attack that undermines their intellect and character, which is a common tactic in sexist hate speech. All individuals, regardless of gender, deserve to have their views and abilities respected without being subjected to such harmful generalizations.} \\
Hybrid RAG & Reducing women to stereotypes ignores the real, diverse experiences and resilience of women, whose voices and skills deserve to be seen as powerful and intelligent. We must actively showcase empowered female role models and challenge narratives that aim to humiliate or objectify women. \\
Hybrid + pred. label & Reducing women to insults ignores the reality that sexist hate speech aims to devalue women's opinions and silence their voices, which has severe psychological and emotional impacts. Every woman deserves to express her views without facing attacks based on her gender. \\
Hybrid + gold label & Labeling women with derogatory terms like 'emotional' or 'empty head' is a form of descriptive attack that undermines their value and silences their voices. Women are capable of being passionate and assertive, and their opinions deserve to be heard without being subjected to sexist objectification. \\
\bottomrule
\end{tabularx}
\end{table*}
